\chapter{Modélisation des prêts hypothécaires}\label{ch:rc:mortgage-modelling}

Le 13 juin 2003, le Treasury à dix ans rapportait 3.13~\%~; le 2 septembre, 4.61~\%. Une partie
de ce mouvement tenait à la couverture que le marché hypothécaire mettait en œuvre pour lui-même. Les propriétaires
américains peuvent rembourser leurs prêts à taux fixe à tout moment~: quand les taux baissent, ils refinancent, quand les taux
montent, ils conservent leur prêt, de sorte qu'un titre adossé à des créances hypothécaires se raccourcit quand les taux baissent
et s'allonge quand ils montent. Les détenteurs qui veulent maintenir la duration de leurs portefeuilles doivent recevoir du taux fixe
quand les taux baissent et en payer quand ils montent~; à l'été 2003, ils ont payé sur un marché qui vendait déjà. Une étude
de la Réserve fédérale publiée cette année-là a montré que la volatilité des swaps
augmentait lorsque le risque de remboursement anticipé était élevé, ce qui est cohérent avec une couverture amplifiant
les mouvements de taux pendant des mois. Le livre~2, chapitre~12, a décrit le marché, ses conventions de remboursement
anticipé et sa convexité négative sur un seul chemin de taux. Ce chapitre modélise le remboursement anticipé sur de nombreux chemins,
valorise un pool par spread ajusté des options, le scinde en sa part d'intérêts et sa part de principal, et
mesure la couverture qu'exige une hausse de
100~points de base.

\section{Anatomie d'un modèle de remboursement anticipé}

\begin{definition}[Modèle de remboursement anticipé]\label{def:rc:mortgage-modelling:model}
Un \emph{modèle de remboursement anticipé}\index{modèle de remboursement anticipé} donne le taux de remboursement anticipé conditionnel
d'un pool chaque mois en fonction des caractéristiques du pool (âge, taux, montant des prêts)
et de l'environnement de taux sur le chemin~; il est estimé sur des données historiques
au niveau de chaque prêt et exécuté chemin par chemin lors de la valorisation.
\end{definition}

\begin{definition}[Rotation du parc immobilier, incitation au refinancement]\label{def:rc:mortgage-modelling:turnover}
La \emph{rotation du parc immobilier}\index{rotation du parc immobilier} (housing turnover) est la part du remboursement anticipé qui provient
des ventes de logements, des défauts et des remboursements partiels, qui croît avec l'âge du prêt jusqu'à un palier et ne dépend pas des taux.
L'\emph{incitation au refinancement}\index{incitation au refinancement} (refinancing incentive) est le taux moyen pondéré du coupon du pool (WAC)
moins le taux hypothécaire qu'un emprunteur peut obtenir aujourd'hui~: le refinancement est rentable
lorsque l'incitation dépasse les coûts qu'il entraîne.
\end{definition}

\begin{definition}[Courbe en S du remboursement anticipé, essoufflement]\label{def:rc:mortgage-modelling:scurve}
La \emph{courbe en S du remboursement anticipé}\index{courbe en S du remboursement anticipé} est la vitesse de refinancement
en fonction de l'incitation~: proche de zéro hors de la monnaie, elle croît fortement sur quelques dizaines de points de base
d'incitation, puis s'aplatit à un maximum. L'\emph{essoufflement}\index{essoufflement} (burnout) est la baisse de la réponse d'un pool à une incitation donnée
après qu'il a déjà été exposé à des incitations~: les emprunteurs qui pouvaient et voulaient refinancer
sont déjà partis, et
ceux qui restent sont plus lents.
\end{definition}

Le modèle de ce chapitre ajoute à la partie rotation, calée sur la rampe PSA
(livre~2, chapitre~12), une courbe en S de refinancement dont le maximum est de 45~\% de CPR et dont le centre se situe
à 75 points de base d'incitation, amortie par l'\omterm{def:rc:mortgage-modelling:scurve}{essoufflement} $e^{-0.25\,c}$ où $c$ désigne l'incitation positive cumulée,
exprimée en années-points de pourcentage
(\cref{fig:rc:mortgage-modelling:scurve}).

\begin{definition}[Coupon courant, spread primaire--secondaire]\label{def:rc:mortgage-modelling:cc}
Le \emph{coupon courant}\index{coupon courant} est le rendement d'un titre hypothécaire à annonce (to-be-announced)
évalué au pair, interpolé entre les coupons qui se traitent. Le \emph{spread primaire--secondaire}\index{spread primaire--secondaire}
est le taux hypothécaire proposé aux emprunteurs (le marché primaire) moins le coupon courant (le marché secondaire)~:
la marge des prêteurs et des garants, qui s'élargit lorsque la capacité de production de crédits est rare
(en 2020, elle s'est située entre 73 et 81 points de base au-dessus de ce que la demande expliquait,
de mars à septembre). Un \omterm{def:rc:mortgage-modelling:model}{modèle de remboursement anticipé} fait correspondre les chemins de taux aux taux hypothécaires des emprunteurs
par ces deux canaux~; le chapitre utilise le taux de swap au pair à dix ans augmenté d'un spread fixe de 175
points de base.
\end{definition}

\omcode{code/firm/mbsoas/firm_mbsoas.py}{36}{49}{Le modèle de remboursement anticipé~: rotation plus courbe en S de refinancement amortie par l'essoufflement.}\label{lst:rc:mortgage-modelling:model}

\begin{omfigure}
\begin{tikzpicture}
\begin{axis}[width=11cm,height=5cm,
  xlabel={incitation au refinancement~: WAC moins taux hypothécaire (points de pourcentage)},ylabel={CPR de refinancement (\%)},
  tick label style={font=\small},label style={font=\small},
  xmin=-2,xmax=3,ymin=0,ymax=50,grid=major,grid style={gray!25},
  table/col sep=comma]
\addplot[omThm,very thick] table[x=inc,y=cpr]{figdata/rates-credit-risk/12-mortgage-modelling/scurve.csv};
\end{axis}
\end{tikzpicture}

\omcaption{La courbe en S de refinancement du \omterm{def:rc:mortgage-modelling:model}{modèle de remboursement anticipé} du chapitre, avant \omterm{def:rc:mortgage-modelling:scurve}{essoufflement}~:
vitesse annuelle de refinancement en fonction de l'incitation. La rotation ajoute quelques pour cent à chaque incitation.
Paramètres illustratifs~; voir le
tutoriel du chapitre.}\label{fig:rc:mortgage-modelling:scurve}
\end{omfigure}

\section{Chemins de taux et dépendance au chemin}

L'\omterm{def:rc:mortgage-modelling:scurve}{essoufflement} fait dépendre le remboursement de chaque mois du chemin suivi jusqu'alors, et le solde restant à rembourser
dépend de l'ensemble des remboursements antérieurs~: un pool ne peut pas être valorisé sur un arbre de taux,
seulement par simulation de chemins complets.

\begin{method}[Valorisation chemin par chemin d'un pool]\label{met:rc:mortgage-modelling:paths}
Simuler des chemins mensuels de taux court avec le \omterm{def:rc:short-rate-models:hw}{modèle de Hull--White} du chapitre~7
(paires antithétiques). Sur chaque chemin, chaque mois~: calculer le taux au pair à dix ans à partir des prix d'obligations du modèle,
le taux hypothécaire, l'incitation et le CPR~; exécuter l'amortissement du pool (mensualité constante, principal programmé,
remboursement anticipé du reste)~; enregistrer les intérêts et le principal versés aux investisseurs.
Actualiser les flux de chaque chemin à son taux court augmenté d'un spread,
puis faire la moyenne.
\end{method}

\omcode{code/firm/mbsoas/firm_mbsoas.py}{74}{94}{Les flux du pool le long de tous les chemins à la fois.}\label{lst:rc:mortgage-modelling:flows}

\section{Le spread ajusté des options en pratique}

Le spread ajusté des options du livre~2, chapitre~21, est le spread constant
ajouté au taux court du modèle qui rend la moyenne des flux actualisés égale
au prix de marché.

\begin{definition}[Coût de l'option]\label{def:rc:mortgage-modelling:optioncost}
Le \emph{coût de l'option}\index{coût de l'option} d'un titre adossé à des créances hypothécaires est son spread
à volatilité nulle (le spread qui le réévalue sur le chemin unique des taux forward) moins son spread ajusté des options~:
la valeur, exprimée en spread, des options de remboursement anticipé des emprunteurs
que la volatilité des taux rend précieuses.
\end{definition}

\begin{example}[Un pool à prime au pair]\label{ex:rc:mortgage-modelling:oas}
Un pool dont le WAC est de 6.0~\%, un coupon de 5.5~\%, 348 mois restants, sur la courbe SOFR du chapitre~1
(taux au pair à dix ans de 3.80~\% selon la convention annuelle du modèle, soit un taux hypothécaire de 5.55~\% et une incitation de 0.45~\%), avec
Hull--White $\kappa = 3~\%$, $\sigma = 90$ points de base et 4\,000 chemins, valorisé à 100~: son spread ajusté des options
est de 146 points de base, son spread à volatilité nulle de 181, et son
\omterm{def:rc:mortgage-modelling:optioncost}{coût de l'option} de 34. Sa duration effective est de 4.99 ans, sa convexité effective $-123$, sa
durée de vie moyenne pondérée de 6.4 ans. Sans \omterm{def:rc:mortgage-modelling:scurve}{essoufflement}, le même pool au même spread vaudrait 99.20~:
un remboursement anticipé plus rapide d'un pool à prime rend le principal au pair plus tôt.
\end{example}

\begin{definition}[Strips intérêts seuls et principal seul]\label{def:rc:mortgage-modelling:iopo}
Un \emph{strip intérêts seuls}\index{strip intérêts seuls} (IO) reçoit les versements d'intérêts d'un pool
et un \emph{strip principal seul}\index{strip principal seul} (PO) ses versements de principal,
programmés et anticipés~; ensemble, ils forment le pool.
\end{definition}

\begin{example}[Découpage du pool]\label{ex:rc:mortgage-modelling:iopo}
Au spread ajusté des options, le pool se découpe en un IO valant 25.26 et un PO valant
74.74. Lorsque les taux montent de 25 points de base, l'IO gagne (26.61) et le PO perd (72.10)~:
un remboursement anticipé plus lent allonge le flux d'intérêts et retarde le principal. L'IO est un
actif rare à duration négative (la \cref{fig:rc:mortgage-modelling:iopo}).
\end{example}

\begin{omfigure}
\begin{tikzpicture}
\begin{axis}[width=11.5cm,height=5.2cm,
  xlabel={décalage parallèle de la courbe (pb)},ylabel={valeur (pour 100 de pool)},
  tick label style={font=\small},label style={font=\small},
  xmin=-210,xmax=210,ymin=15,ymax=90,grid=major,grid style={gray!25},
  legend columns=2,legend style={font=\footnotesize,at={(0.5,-0.25)},anchor=north,draw=none,fill=none,/tikz/every even column/.append style={column sep=8pt}},
  table/col sep=comma]
\addplot[omThm,very thick,mark=*] table[x=shift,y=po]{figdata/rates-credit-risk/12-mortgage-modelling/iopo.csv};
\addplot[omDef,very thick,mark=square*] table[x=shift,y=io]{figdata/rates-credit-risk/12-mortgage-modelling/iopo.csv};
\legend{strip principal seul,strip intérêts seuls}
\end{axis}
\end{tikzpicture}

\omcaption{Valeurs des \omterm{def:rc:mortgage-modelling:iopo}{strips principal seul} et intérêts seuls du pool lorsque la courbe entière
se déplace, au spread ajusté des options du pool. Le PO monte fortement lorsque les taux baissent
et que les remboursements anticipés s'accélèrent~; l'IO monte avec les taux. Données~: tutoriel du
chapitre.}\label{fig:rc:mortgage-modelling:iopo}
\end{omfigure}

\section{Couverture de la convexité négative}

\begin{omfigure}
\begin{tikzpicture}
\begin{axis}[width=11.5cm,height=5.2cm,
  xlabel={décalage parallèle de la courbe (pb)},ylabel={valeur du pool},
  tick label style={font=\small},label style={font=\small},
  xmin=-210,xmax=210,ymin=87,ymax=112,grid=major,grid style={gray!25},
  legend columns=2,legend style={font=\footnotesize,at={(0.5,-0.25)},anchor=north,draw=none,fill=none,/tikz/every even column/.append style={column sep=8pt}},
  table/col sep=comma]
\addplot[omThm,very thick,mark=*,mark size=1.5pt] table[x=shift,y=value]{figdata/rates-credit-risk/12-mortgage-modelling/price.csv};
\addplot[omDef,thick,dashed] table[x=shift,y=linear]{figdata/rates-credit-risk/12-mortgage-modelling/price.csv};
\legend{valeur du pool (OAS maintenu),tangente à la courbe d'aujourd'hui}
\end{axis}
\end{tikzpicture}

\omcaption{Valeur du pool en fonction d'un décalage parallèle de la courbe~: en dessous de sa tangente des deux côtés,
signature de la convexité négative. Ses gains lorsque les taux baissent sont plafonnés par le refinancement~;
ses pertes lorsque les taux montent croissent à mesure que sa duration s'allonge. Données~: tutoriel du
chapitre.}\label{fig:rc:mortgage-modelling:price}
\end{omfigure}

Un détenteur couvert à DV01 nul est vendeur de gamma~: après une hausse, il redevient long de duration
et doit payer le fixe (ou vendre des bons du Trésor) pour se recouvrir~; après une baisse, il doit recevoir le
fixe. Ses transactions de couverture vont dans le sens du mouvement, et lorsque de nombreux détenteurs
appliquent la même couverture, ces transactions font bouger le marché.

\begin{example}[Cent points de base]\label{ex:rc:mortgage-modelling:event}
Après une hausse de 100 points de base, au même spread ajusté des options, le pool vaut 94.50
et sa duration s'est allongée de 4.99 à 6.17 ans. Sur 10~milliards de dollars de nominal, le
DV01 passe de 4.99~millions de dollars à 5.83 millions par point de base. Un swap au pair à dix ans sur
la courbe décalée a un DV01 de 796\,282 dollars par milliard, si bien que le détenteur doit payer le fixe sur
1.06~milliard de dollars de swaps à dix ans pour être de nouveau couvert.
\end{example}

Qui réalise ces couvertures compte autant que leur ampleur. Un détenteur qui ne couvre pas la duration,
comme une banque centrale qui achète pour des raisons de politique monétaire, absorbe la convexité négative
sans la reporter sur le marché~; lorsque ce détenteur cesse d'acheter et laisse ses avoirs s'éteindre,
la convexité revient aux mains privées, qui la couvrent.

\begin{omfigure}
\begin{tikzpicture}
\begin{axis}[width=11.5cm,height=5cm,
  xlabel={année},ylabel={billions de dollars américains},
  tick label style={font=\small},label style={font=\small},
  xmin=2003,xmax=2027,ymin=0,ymax=3,grid=major,grid style={gray!25},
  xtick={2005,2010,2015,2020,2025},xticklabels={2005,2010,2015,2020,2025},
  table/col sep=comma]
\addplot[omThm,very thick] table[x=year,y=tn]{figdata/rates-credit-risk/12-mortgage-modelling/fedmbs.csv};
\end{axis}
\end{tikzpicture}

\omcaption{Titres adossés à des créances hypothécaires détenus en propre par la Réserve fédérale, niveaux de fin de mois
à un mercredi, de décembre 2002 à août 2026~: nuls avant 2009, un sommet de 2.72 billions de dollars américains en février
2022, 1.91 billion de dollars américains en août 2026. Source~: Federal Reserve H.4.1
via FRED (série WSHOMCB, fichier de données du livre~2).}\label{fig:rc:mortgage-modelling:fed}
\end{omfigure}

\begin{dated}{2026-09}{Qui détient la convexité}\label{dat:rc:mortgage-modelling:holders}
Les avoirs de la Réserve fédérale en titres adossés à des créances hypothécaires d'agence ont culminé à 2.72 billions
de dollars américains en février 2022 et s'établissaient à 1.91 billion de dollars américains en août 2026, en décroissance
à mesure que les pools se remboursent par anticipation. Les titres hypothécaires d'agence se sont échangés à environ
367 milliards de dollars par jour en 2026 jusqu'en août, principalement sur le marché du to-be-announced (SIFMA).
\end{dated}

\section{Risque de modèle dans le remboursement anticipé}

Chaque chiffre ci-dessus repose sur un modèle comportemental~: le centre et la hauteur de la courbe en S,
le taux d'\omterm{def:rc:mortgage-modelling:scurve}{essoufflement}, l'écart de taux hypothécaire. Ils sont estimés sur des vagues de refinancement passées,
et chaque nouvelle vague diffère (technologies qui accélèrent le refinancement, critères de souscription plus stricts
qui le ralentissent, un \omterm{def:rc:mortgage-modelling:cc}{spread primaire--secondaire} qui s'élargit lorsque les établissements prêteurs sont saturés).
Les desks rapportent donc le spread ajusté des options selon plusieurs modèles,
calculent des «~durations~» de remboursement anticipé (la sensibilité à un changement d'échelle des vitesses du modèle),
et constituent des réserves pour l'incertitude de modèle (chapitres~26 et~27).

\section{Tutoriel~: un pool sur de nombreux chemins}\label{tut:rc:mortgage-modelling:tutorial}

\begin{tutorial}
\textbf{Objectif.} Évaluer un pass-through au spread ajusté des options, le décomposer en strips
et mesurer sa convexité négative. \textbf{État final~:}
la \cref{fig:rc:mortgage-modelling:iopo,fig:rc:mortgage-modelling:price} et les résultats des
\cref{ex:rc:mortgage-modelling:oas,ex:rc:mortgage-modelling:event}.
\begin{tutsteps}
\item \textbf{Trajectoires}~: \texttt{simulate\_paths(curve, kappa, sigma, months, paths)} renvoie
les taux courts mensuels et les taux au pair à dix ans.
\item \textbf{Flux}~: \texttt{pool\_flows(pool, model, par10)}.
\item \textbf{Spread et risque}~: \texttt{base()}, \texttt{duration\_at(0.01, oas)}.
\item \textbf{Couverture}~: \texttt{convexity\_event()}~; \texttt{fig\_rc\_mbs.py} produit les
graphiques.
\end{tutsteps}
\textbf{Pour aller plus loin.} Doubler la pente de la courbe en S et comparer la convexité~;
relever le spread hypothécaire de 50 points de base (un \omterm{def:rc:mortgage-modelling:cc}{spread primaire--secondaire} élargi) et
mesurer la variation de valeur.
\end{tutorial}

\section{Construction~: le moteur de spread ajusté des options}\label{bld:rc:mortgage-modelling:mbsoas}

\begin{build}
\bfield{Rôle} Le portefeuille hypothécaire de l'entreprise miniature~: remboursement anticipé
par trajectoire, spreads ajustés des options, strips et risque effectif.
\bfield{Interface} \texttt{Pool}, \texttt{PrepayModel} (avec \texttt{refi},
\texttt{cpr})~; \texttt{simulate\_paths}~; \texttt{pool\_flows}~; \texttt{pv\_paths}~;
\texttt{solve\_oas}~; \texttt{analyse(curve, pool, model, price, kappa, sigma, paths)}.
\bfield{Règles} Pas mensuels~; le \texttt{firm\_prepay} du livre~2 pour l'amortissement et la
rampe PSA~; nombres aléatoires communs pour toute sensibilité~; numpy uniquement.
\bfield{Tests d'acceptation} \texttt{code/firm/mbsoas/tests/}~: sans remboursement anticipé, le
pool est un prêt amortissable sur chaque trajectoire~; le spread ajusté des options se reprice et les
strips se somment~; un pool proche de la monnaie a une convexité négative et un coût d'option positif, son IO
gagne et son PO perd lorsque les taux montent~; l'\omterm{def:rc:mortgage-modelling:scurve}{essoufflement} ralentit le remboursement anticipé.
\bfield{Pour aller plus loin} Un modèle au niveau du prêt à partir de données de cohortes~; un
modèle de taux à deux facteurs~; le financement implicite du dollar roll sur les flux du modèle.
\end{build}

\begin{omsources}
\item R.\ Perli et B.\ Sack, «~Does mortgage hedging amplify movements in long-term
interest rates?~», Finance and Economics Discussion Series 2003-49, Federal Reserve Board.
\item US Department of the Treasury, \emph{Daily Treasury Par Yield Curve Rates}, 2003.
\item L.\ Hayre (dir.), \emph{Salomon Smith Barney Guide to Mortgage-Backed and
Asset-Backed Securities}, Wiley, 2001.
\item A.\ Fuster, A.\ Hizmo, L.\ Lambie-Hanson, J.\ Vickery et P.\ S.\ Willen, «~How resilient is
mortgage credit supply? Evidence from the COVID-19 pandemic~», NBER Working Paper 28843.
\end{omsources}

\section{Exercices}

\begin{exercise}[$\star$]\label{exo:rc:mortgage-modelling:1}
Un pool a un WAC de 6.0~\% et le taux au pair à dix ans est de 3.80~\%. Avec un spread hypothécaire
de 1.75~\%, quelle est l'\omterm{def:rc:mortgage-modelling:turnover}{incitation au refinancement}~?
\end{exercise}

\begin{exercise}[$\star$]\label{exo:rc:mortgage-modelling:2}
En utilisant l'\cref{ex:rc:mortgage-modelling:oas}, quel est le \omterm{def:rc:mortgage-modelling:optioncost}{coût de l'option} et que
mesure-t-il~?
\end{exercise}

\begin{exercise}[$\star$]\label{exo:rc:mortgage-modelling:3}
Pourquoi l'IO gagne-t-il lorsque les taux montent~?
\end{exercise}

\begin{exercise}[$\star\star$]\label{exo:rc:mortgage-modelling:4}
Pourquoi un pool ne peut-il pas être évalué sur un arbre recombinant, contrairement à une option bermudéenne~?
\end{exercise}

\begin{exercise}[$\star\star$]\label{exo:rc:mortgage-modelling:5}
Lire la \cref{fig:rc:mortgage-modelling:price}~: estimer le gain lié à une baisse de 100 points de base
et la perte liée à une hausse de 100 points de base, puis les comparer.
\end{exercise}

\begin{exercise}[$\star\star$]\label{exo:rc:mortgage-modelling:6}
Sans \omterm{def:rc:mortgage-modelling:scurve}{essoufflement}, le pool vaut 99.20 au même spread. Expliquer le signe.
\end{exercise}

\begin{exercise}[$\star\star\star$]\label{exo:rc:mortgage-modelling:7}
\emph{Programmation.} Reproduire la reconstitution de la couverture de l'\cref{ex:rc:mortgage-modelling:event}~: le DV01
avant et après la hausse sur 10~milliards de dollars, et le notionnel du swap à dix ans.
\end{exercise}

\begin{exercise}[$\star\star\star$]\label{exo:rc:mortgage-modelling:8}
\emph{Trouver la faille.} «~Notre book de MBS est couvert en duration chaque jour, il n'a donc aucun risque de taux.~»
\end{exercise}

\section{Problème~: l'événement de convexité}

\begin{problem}[{Problème du week-end --- un portefeuille rétablit sa couverture après une chute des cours}]\label{pb:rc:mortgage-modelling:1}
Un investisseur hypothécaire détient 10~milliards de dollars du pool du chapitre, couvert à un DV01 nul
au moyen de swaps payeurs à dix ans. Les taux montent de 100 points de base en parallèle.

\medskip\noindent\textbf{Partie I --- Avant.}
\begin{enumerate}
\item Donner le spread ajusté des options, le spread à volatilité nulle et le \omterm{def:rc:mortgage-modelling:optioncost}{coût de l'option} du pool.
\item Donner sa duration et sa convexité.
\item Donner le DV01 de la position, par point de base.
\item Quelle part de la valeur provient des intérêts, quelle part du principal~?
\item Quelle position de l'investisseur est vendeuse d'options, et envers qui~?
\end{enumerate}

\medskip\noindent\textbf{Partie II --- La hausse.}
\begin{enumerate}[resume]
\item Donner la valeur du pool après la hausse, au même spread.
\item Donner sa nouvelle duration et le nouveau DV01 de la position.
\item Pourquoi la duration s'est-elle allongée~?
\item Donner le notionnel de swaps payeurs à dix ans nécessaire pour rétablir la couverture.
\item Que se passe-t-il si chaque investisseur hypothécaire fait la même opération le même jour~?
\end{enumerate}

\medskip\noindent\textbf{Partie III --- Strips.}
\begin{enumerate}[resume]
\item Donner les valeurs de l'IO et du PO pour $\pm25$ points de base.
\item Quel strip couvre la convexité de l'autre~?
\item Pourquoi les IO attirent-ils les banques dont les dépôts gagnent en valeur lorsque les taux montent~?
\item Que perd le détenteur d'un PO si le remboursement anticipé ralentit pour des raisons étrangères aux taux~?
\item Quel strip est le plus exposé au risque de \omterm{def:rc:mortgage-modelling:model}{modèle de remboursement anticipé}~?
\end{enumerate}

\medskip\noindent\textbf{Partie IV --- Jugement.}
\begin{enumerate}[resume]
\item Faut-il couvrir avec des swaptions plutôt qu'avec des swaps~? Combien cela coûterait-il~?
\item Pourquoi l'épisode de 2003 a-t-il duré des mois plutôt que des jours~?
\item Comment une banque centrale détenant des titres hypothécaires modifierait-elle la dynamique~?
\item Énoncer le \emph{résultat nommé}~: le DV01 avant et après, ainsi que le notionnel du swap de
recouverture (re-hedge).
\item En une phrase~: pourquoi une hypothèque s'allonge-t-elle au moment où l'on en voudrait le moins~?
\end{enumerate}
\end{problem}

\section{Questions d'entretien}

\begin{interviewq}[$\star$ \iqroles{trader, chercheur}]\label{iq:rc:mortgage-modelling:1}
Pourquoi les titres adossés à des hypothèques ont-ils une convexité négative~?
\end{interviewq}

\begin{interviewq}[$\star\star$ \iqroles{chercheur}]\label{iq:rc:mortgage-modelling:2}
Que contient un \omterm{def:rc:mortgage-modelling:model}{modèle de remboursement anticipé}, et quelles composantes dépendent des taux~?
\end{interviewq}

\begin{interviewq}[$\star\star$ \iqroles{chercheur, développeur}]\label{iq:rc:mortgage-modelling:3}
Comment calculez-vous un spread ajusté des options, et pourquoi recourir à Monte Carlo~?
\end{interviewq}

\begin{interviewq}[$\star\star$ \iqroles{trader}]\label{iq:rc:mortgage-modelling:4}
Les taux remontent de 50 points de base. Que font les hedgers hypothécaires, et quel effet cela a-t-il sur le
marché~?
\end{interviewq}

\begin{interviewq}[$\star\star\star$ \iqroles{risque, chercheur}]\label{iq:rc:mortgage-modelling:5}
Comment mesureriez-vous le risque de \omterm{def:rc:mortgage-modelling:model}{modèle de remboursement anticipé}, et quelles réserves constitueriez-vous pour ce risque~?
\end{interviewq}

\begin{interviewq}[$\star\star\star$ \iqroles{développeur}]\label{iq:rc:mortgage-modelling:6}
Votre duration effective saute d'un jour à l'autre de six mois, sans mouvement de marché. Que
vérifiez-vous~?
\end{interviewq}
